%!LW recipe=pLaTeX (ptex2pdf) -> pbibtex -> ptex2pdf*2

%\documentclass[submit,techrep]{ipsj}
\documentclass[submit,techrep,noauthor]{ipsj}



\usepackage[dvips]{graphicx}
\usepackage{latexsym}

\def\Underline{\setbox0\hbox\bgroup\let\\\endUnderline}
\def\endUnderline{\vphantom{y}\egroup\smash{\underline{\box0}}\\}
\def\|{\verb|}


\makeatletter
\providecommand{\newblock}{\hskip .11em\@plus.33em\@minus.07em}
\makeatother
%

%\setcounter{巻数}{59}%vol59=2018
%\setcounter{号数}{10}
%\setcounter{page}{1}


\begin{document}


\title{サンプリングと階層化による\\
$k$-近傍グラフマージアルゴリズム}

\etitle{A $k$-NN Graph Merging Algorithm\\
 on Sampling and Hierachical Structure}

\affiliate{ISCT}{東京科学大学\\
2--12--1, Ookayama, Meguro-ku, Tokyo 152--8552, Japan}


\author{橋本 龍徳}{Hashimoto Tatsunori}{ISCT}[hashimoto@lsc.comp.isct.ac.jp]
\author{ベルジャブ ネスリーン}{Berjab Nesrine}{ISCT}[berjab@comp.isct.ac.jp]
\author{常 穹}{Chang Qiong}{ISCT}[q.chang@c.titech.ac.jp]
\author{宮崎 純}{Miyazaki Jun}{ISCT}[miyazaki@comp.isct.ac.jp]

\begin{abstract}
大規模データやデータの逐次的な増加に対応する昨今、独立に構築された$k$-近傍（$k$-NN）グラフをマージすることは有用である。しかし、既存手法では近傍点同士の比較に基づく精製のコストが大きく、近傍数$k$の増加に伴い処理時間が増大するという課題がある。本研究では、サブグラフ内の既存の近傍関係を活用しつつ、部分グラフ間の近傍探索に重点を置いた二層型$k$-NNグラフマージ手法を提案する。提案手法では、サンプリングした点からなる小規模な上位層を用いて探索の開始点を求め、$k$に依存しない探索幅でグラフを探索することで近傍候補を見つける。さらに、対称性の活用やPCA変換空間における距離計算の早期打ち切りにより計算コストを削減する。複数のデータセットを用いた実験により、提案手法は既存手法と比較して$k$の増加に伴う処理時間の増大を抑制し、高速なグラフマージを実現できることを示す。
\end{abstract}


%
%\begin{jkeyword}
%情報処理学会論文誌ジャーナル，\LaTeX，スタイルファイル，べからず集
%\end{jkeyword}
%
%\begin{eabstract}
%This document is a guide to prepare a draft for submitting to IPSJ
%Journal, and the final camera-ready manuscript of a paper to appear in
%IPSJ Journal, using {\LaTeX} and special style files.  Since this
%document itself is produced with the style files, it will help you to
%refer its source file which is distributed with the style files.
%\end{eabstract}
%
%\begin{ekeyword}
%IPSJ Journal, \LaTeX, style files, ``Dos and Dont's'' list
%\end{ekeyword}

\maketitle

%1
\section{はじめに}
$k$-近傍グラフは、各データ点をその点に最も近い$k$個のデータ点と接続するグラフであり、データ点間の局所的な近傍関係を表現するための基本的なデータ構造である。
$k$-NNグラフは、機械学習 \cite{boiman_defense_2008,flores_analysis_2022}、多様体学習 \cite{yan_graph_2007,belkin_laplacian_2003,roweis_nonlinear_2000}、情報検索 \cite{li_approximate_2020}、データマイニング \cite{brito_connectivity_1997} など、幅広い分野で利用されている。

さらに、$k$-NNグラフは、グラフベースの近似最近傍探索（Approximate Nearest Neighbor Search; ANNS）インデックスを構築するための基盤にもなる。例えば、NSG \cite{fu_fast_2019} や CAGRA \cite{ootomo_cagra_2024} は、初期の近似 $k$-NNグラフを改良したり、枝刈りを行ったりすることで、効率的な探索グラフを構築する。

したがって、効率的な $k$-NNグラフの構築は、近傍関係を直接利用するアプリケーションだけでなく、大規模な類似検索においても重要である。
2. $k$-NNグラフ構築の計算コスト
\(n\) 個の \(d\) 次元データ点に対して、全探索によって厳密な $k$-NNグラフを構築するには、\(O(dn^2)\) の時間計算量を要するため、大規模なデータセットに対しては現実的ではない。

この問題は、特に大規模かつ高次元のデータセットにおいて顕著になる。

NN-Descent \cite{dong_efficient_2011} や、そのGPU実装であるGNND \cite{wang_fast_2021} などの近似構築アルゴリズムは、計算コストを大幅に削減する。

しかし、これらの手法を用いた場合でも、大規模なデータセット全体に対して一度にグラフを構築することは、依然として困難な場合がある。
3. グラフマージによる構築
これに対する別のアプローチとして、マージに基づく構築手法がある。

この手法では、データセットをより小さな部分集合に分割し、それぞれの部分集合に対して独立に $k$-NNグラフを構築した後、それらをマージすることで一つの大規模なグラフを得る。

このアプローチにより、サブグラフを独立に構築し、必要に応じて異なる計算デバイス上で処理した後に、それらをマージすることが可能になる。
4. 既存グラフの再利用と逐次的な更新
さらに重要なこととして、グラフマージは、サブグラフがすでに存在する場合にも有用である。

例えば、新たなデータが追加された際には、追加されたデータ点に対して $k$-NNグラフを構築し、それを既存のグラフとマージすることができる。

このように、サブグラフ内の近傍関係を再利用することで、グラフを最初から構築し直す必要がなくなる。

より一般的には、独立に構築されたグラフが持つ既存の近傍情報を破棄することなく、それらをマージすることが可能である。

これらの性質から、効率的なグラフマージは、$k$-NNグラフのスケーラブルな構築と逐次的な維持・更新の両方において有用である。
5. 既存手法とその課題
しかし、既存のグラフマージ手法は、NN-Descentに基づく改良処理に大きく依存している。

S-Merge \cite{zhao_merge_2022} は、ランダムに選択した候補点を用いてサブグラフ間の接続を初期化する。一方、LG²-Merge \cite{hashimoto_lg2merge} は、局所性鋭敏型ハッシュ（Locality-Sensitive Hashing; LSH）を利用して、より有望な候補点を選択することで、その後の改良処理のコストを削減する。

しかし、これらの手法は初期化処理が改善されているものの、いずれもNN-Descentに基づく改良処理に依存しており、各データ点について、各反復で最大 \(O(k^2)\) 回の候補点の比較が必要となる場合がある。

その結果、$k$ が大きくなるにつれて、改良処理のコストが急速に増大する。

これは、ANNSインデックスのために比較的大きな $k$ を持つ初期グラフを構築する場合に、特に問題となる。
6. 提案手法 TL-Merge
本論文では、NN-Descentに基づく改良処理を必要とせず、サブグラフ間の近傍点を効率的に発見する手法、Two-Layer $k$-NN Graph Merging（TL-Merge）を提案する。

本手法の中心となるアイデアは、入力グラフに既存の接続関係を再利用し、サブグラフをまたぐ近傍点の探索を誘導することである。

ここでは、二つのサブグラフをマージする場合を例に、その処理を説明する。

TL-Mergeは、まず各サブグラフから少数の固定数のデータ点をサンプリングし、コンパクトな上位層を構築する。上位層はサンプリングされたデータ点のみから構成され、それらの間に辺は存在しない。

一方のサブグラフに含まれる各データ点に対して、TL-Mergeはもう一方のサブグラフの上位層を利用して、有望な探索開始点を特定する。そして、その探索開始点から既存の近傍リストをたどることで、近傍点を発見する。

発見されたサブグラフ間の近傍候補は、その後、近傍リストに挿入される。

以降の反復処理では、新たに発見されたサブグラフ間の近傍点を含む、更新済みの近傍リストを利用して、さらに候補点を探索し、グラフの精度を段階的に向上させる。

さらに、この改良処理において、各データ点・各反復あたりの候補点に対する距離計算回数は、二つの探索幅パラメータの積によって上限が定められる。これらのパラメータは、$k$ とは独立に設定することができる。

加えて、TL-Mergeは、PCAに基づく距離計算の早期打ち切りと、双方向の近傍更新を利用することで、さらなる効率化を実現する。
7. 本論文の主な貢献
本論文の主な貢献は、以下のとおりである。

1. TL-Mergeの提案

サブグラフ内の既存の近傍関係を再利用し、サブグラフ間の近傍点の発見に重点を置く、二層構造のグラフマージ手法TL-Mergeを提案する。探索幅を固定したグラフウォークにより、候補点に対する距離計算回数が $k$ に対して二次的に増大することを回避する。

2. 距離計算の効率化

   冗長な計算を削減するため、PCAに基づく距離計算の早期打ち切りと、双方向の近傍更新を導入する。前者は、距離計算の対象となる二つのデータ点のいずれの近傍リストも改善できない場合に、その距離計算を安全に打ち切る。一方、後者は、距離計算の対称性を利用する。

3. 既存手法に対する高速化とスケーラビリティ

五つのデータセットを用いた実験により、TL-Mergeは \(k=100\) において、S-Mergeに対して最大8.1倍の高速化を達成することを示す。

さらに、TL-Mergeは $k$ に対して優れたスケーラビリティを示す。具体的には、$k$ を20から100に増加させた際、既存手法ではマージ時間が13～24倍に増加するのに対して、TL-Mergeでは2.6～4.8倍の増加にとどまる。

4. グラフの精度とパラメータに関する評価

   さらに、五つすべてのデータセットにおいて、TL-Mergeは比較手法よりも短い時間で同一のRecall目標値に到達し、近傍距離の誤差も小さいことを示す。

   また、パラメータ感度に関する実験から、比較的小規模な上位層を使用し、網羅的なパラメータ調整を行わなくても、良好な性能を達成できることを示す。
8. 論文の構成
本論文の残りの構成は、以下のとおりである。

第 \ref{sec:related-work} 節では、グラフマージ問題を定義し、関連研究を概観する。第 \ref{sec:proposed} 節ではTL-Mergeについて説明し、第 \ref{sec:experiments} 節ではその性能を評価する。最後に、第 \ref{sec:conc} 節で本論文をまとめ、今後の課題について述べる。




%2
\section{投稿の流れ}


%2.1
\subsection{準備}

情報処理学会論文誌ジャーナルの\LaTeX スタイルファイルを含む論文執筆キットは
\begin{quote}
aaa
\end{quote}
からダウンロードすることができる．論文執筆キットは以下のファイルを含んで
いる．
実行環境としては\LaTeXe を前提としているので，準備されたい．







%2.2
\subsection{原稿の作成と投稿}

本稿に従って用意した投稿用原稿の\LaTeX ソースからpdfファイルを作成し，
(**部分は研究会の略称，DBS等)の研究会投稿システムにて，指示にし従い投稿する．







%3
\section{論文フォーマットの指針}
\label{sec:format}

以下，情報処理学会論文誌ジャーナル用スタイルファイルを用いた論文フォーマットの指針について述べるので，
これに従って原稿を用意頂きたい．
\LaTeX を用いた一般的な文章作成技術については，
\cite{okumura, companion} 等を参考にされたい．



%4
\section{論文の構成}
\label{config}

ファイルは次のようになる．下線部は投稿時に省略可能なもの．


%4.1
\subsection{表題・著者名等}

表題，著者名とその所属，および概要を前述のコマンドや環境により{\bf 和文と
英文の双方について}定義した後， によって出力する．


\newpage%%%%%

%4.1.1
\subsubsection{表題} 

%4.2.3
\subsubsection{フォントサイズ}

フォントサイズは，スタイルファイルによって自動的に設定されるため，
基本的には著者が自分でフォントサイズを変更する必要はない．

%4.2.4
\subsubsection{句読点}

句点には全角の「．」，
読点には全角の「，」を用いる．
ただし英文中や数式中で「.」や「,」を使う場合には，
半角文字を使う．「。」や「、」は使わない．



%4.2.5
\subsubsection{全角文字と半角文字}

全角文字と半角文字の両方にある文字は次のように使い分ける．

\begin{enumerate}
\item 括弧は全角の「（」と「）」を用いる．但し，英文の概要，図表見出し，
書誌データでは半角の「(」と「)」を用いる．

\item 英数字，空白，記号類は半角文字を用いる．ただし，句読点に関しては，
前項で述べたような例外がある．

\item カタカナは全角文字を用いる．

\end{enumerate}

%4.2.6
\subsubsection{箇条書}


%4.2.7
\subsubsection{脚注}



また場合によっては，
脚注をつけた位置と脚注本体とを別の段に置く方がよいこともある．


なお，脚注番号は論文内で通し番号で出力される．




%4.2.8
\subsubsection{OverfullとUnderfull}

組版時にはoverfullを起こさないことを原則としている．
従って，まず提出するソースが著者の環境でoverfullを起こさないように，
文章を工夫するなどの最善の努力を払っていただきたい．

また著者の執筆時点では発生しないoverfullが，
組版時の環境では発生することもある．
このような事態をできるだけ回避するために，
パラグラフの先頭付近では長い英単語を使用しない，
などの注意を払うようにして頂きたい．




%4.3
\subsection{数式}\label{sec:Item}

%4.3.1
\subsubsection{本文中の数式}

\subsubsection{別組の数式}

%
\begin{equation}
\Delta_l = \sum_{i=l|1}^L\delta_{pi}
\end{equation}
%
のように，センタリングではなく固定字下げで数式を出力し，
かつ背が高い数式による行送りの乱れを吸収する機能がある．

%4.3.3
\subsubsection{eqnarray環境}

互いに関連する別組の数式が2行以上連続して現れる場合には，
．


%4.3.4
\subsubsection{数式のフォント}

\LaTeX が標準的にサポートしているもの以外の特殊な数式用フォントは，
できるだけ使わないようにされたい．
どうしても使用しなければならない場合には，
その旨申し出て頂くとともに，
組版工程に深く関与して頂くこともあることに留意されたい．







%4.4
\subsection{図}

1段の幅におさまる図は，



また紙面スペースの節約のために，




2段の幅にまたがる図は，



図の中身では本文と違い，
どのような大きさのフォントを使用しても構わない（\figref{fig:double} 参照）．
また図の中身として，encapsulate されたPostScriptファイル（いわゆるEPSファイル）を読み込むこともできる．
読み込みのためには，プリアンブルで
%




%4.5
\subsection{表}

表の罫線はなるべく少なくするのが，仕上がりをすっきりさせるコツである．
罫線をつける場合には，
一番上の罫線には二重線を使い，左右の端には縦の罫線をつけない （\tabref{tab:example}）．



また，表の上に和文と英文の双方の見出しを，


\begin{table}[tb] 
\caption{表の例} 
\ecaption{An Example of Table.}
\label{tab:example}
\hbox to\hsize{\hfil
\begin{tabular}{l|lll}\hline\hline
& column1 & column2 & column3 \\\hline
row1 &	item 1,1 & item 2,1 & ---\\
row2 &	---      & item 2,2 & item 3,2 \\
row3 &	item 1,3 & item 2,3 & item 3,3 \\
row4 &	item 1,4 & item 2,4 & item 3,4 \\\hline
\end{tabular}\hfil}
\end{table}




\newpage%%%%%

%4.6
\subsection{参考文献・謝辞}

%4.6.1
\subsubsection{参考文献の参照}

%

が得られる．

%4.6.2
\subsubsection{参考文献リスト}
参考文献リストには，
原則として本文中で引用した文献のみを列挙する．
順序は参照順あるいは第一著者の苗字のアルファベット順とする．
文献リストはBiB\TeX と\verb+ipsjunsrt.bst+（参照順）
または\verb+ipsjsort.bst+（アルファベット順）を用いて作り，
\verb+\bibliograhpystyle+と\verb+\bibliography+コマンドにより
利用することが出来る．
これらを用いれば，
規定の体裁にあったものができるので，
できるだけ利用していただきたい．
また製版用のファイル群には\verb+.bib+ファイルではなく\verb+.bbl+ファイルを
必ず含めることに注意されたい．
一方，何らかの理由でthebibliography環境で文献リストを
「手作り」しなければならない場合は，
このガイドの参考文献リストを注意深く見て，
そのスタイルにしたがっていただきたい．




%4.6.3
\subsubsection{謝辞}





%5
\section{論文内容に関する指針}

論文の内容について，
論文誌ジャーナル編集委員会で作成した「べからず集」を以下に示す．
投稿前のチェックリストとして利用頂きたい．
これ以外にも，査読者用，
メタ査読者用の「べからず集」\cite{webpage2}も公開しているので，
参照されたい．
また，作文技術に関する \cite{book1, book2, book3, book4}のような書籍も参考になる．



%5.1
\subsection{書き方の基本}

\begin{itemize}
 \item[$\Box$] 研究の新規性，有用性，信頼性が読者に伝わるように記述する．
 \item[$\Box$] 読み手に，読みやすい文章を心がける（内容が前後する，背景・
	       課題の設定が不明瞭などは読者にとって負担）．
 \item[$\Box$] 解決すべき問題が汎用化（一般的に記述）されていないのは再
	       考を要する（XX大学の問題という記述に終始）．あるいは，
	       （単に「作りました」だけで）解決すべき問題そのものの記述
	       がないのは再考を要する．
 \item[$\Box$] 結論が明確に記されていない，または，範囲，限界，問題点な
	       どの指摘が適切ではない，または，結論が内容にそったもので
	       はないものは再考を要する．
 \item[$\Box$] 科学技術論文として不適当な表現や，分かりにくい表現がある
	       のは再考を要する．
 \item[$\Box$] 極端な口語体や，長文の連続などは再考を要する．
 \item[$\Box$] 章，節のたて方，全体の構成等が適切でない文章は再考を要す
	       る．
 \item[$\Box$] 文中の文脈から推測しないと内容の把握が困難な論文にしない．
 \item[$\Box$] 説明に飛躍した点があり，仮説等の説明が十分ではないのは再
	       考を要する．
 \item[$\Box$] 説明に冗長な点，逆に簡単すぎる点があるのは再考を要する．
 \item[$\Box$] 未定義語を減らす．
\end{itemize}


%5.2
\subsection{新規性と有効性を明確に示す}

\begin{itemize}
 \item[$\Box$] 在来研究との関連，研究の動機，ねらい等が明確に説明されて
	       いないのは再考を要する．
 \item[$\Box$] 既知／公知の技術が何であって，何を新しいアイデアとして提
	       案しているのかが書かれていないのは再考を要する．
 \item[$\Box$] 十分な参考文献は新規性の主張に欠かせない．
 \item[$\Box$] 提案内容の説明が，概念的または抽象的な水準に終始していて，
	       読者が提案内容を理解できない（それだけで新規性が感じられ
	       ないもの）のは再考を要する．
 \item[$\Box$] 論文で提案した方法の有効性の主張がない，またはきわめて貧
	       弱なのは再考を要する．
\end{itemize}

%5.3
\subsection{書き方に関する具体的な注意}

\begin{itemize}
 \item[$\Box$] 和文標題が内容を適切に表現していないのは再考を要する．
 \item[$\Box$] 英文標題が内容を適切に表現していない，または英語として適
	       切でないのは再考を要する．
 \item[$\Box$] アブストラクトが主旨を適切に表現していない，または英文が
	       適切ではないのは再考を要する．
 \item[$\Box$] 記号・略号等が周知のものでなく，または，用語が適切でなく，
	       または，図・表の説明が適当ではないのは再考を要する．
 \item[$\Box$] 個人的あるいは非常に小さなグループ／企業だけで通用するよ
	       うな用語が特別な説明もなしに多用されているのは再考を要す
	       る．
 \item[$\Box$] 図表自体は十分に明確ではない，または誤りがあるのは再考を
	       要する．
 \item[$\Box$] 図表が鮮明ではないのは再考を要する．
 \item[$\Box$] 図表が大きさ，縮尺の指定が適切でないのは再考を要する.
\end{itemize}

%5.4
\subsection{参考文献}

\begin{itemize}
 \item[$\Box$] 参考文献は10件以上必要（分野によっては20件以上，30件以上
	       という意見もある）.
 \item[$\Box$] 十分な参考文献は新規性の主張に欠かせない．
 \item[$\Box$] 適切な文献が引用されておらず，その数も適切ではないのは再
	       考を要する．
 \item[$\Box$] 日本人によるしかるべき論文を引用することで日本人研究コミュ
	       ニティの発展につながる．
 \item[$\Box$] 参考文献は自分のものばかりではだめ．
\end{itemize}

%5.5
\subsection{二重投稿}

\begin{itemize}
 \item[$\Box$] 二重投稿はしてはならない ─ ただし国際会議に採択された論
	       文を著作権が問題にならないように投稿することは構わない．
 \item[$\Box$] 他の論文とまったく同じ図表を引用の明示なしに利用すること
	       は禁止．
 \item[$\Box$] 既発表の論文等との間に重複があるのは再考を要する．
\end{itemize}

%5.6
\subsection{他の人に読んでもらう}

\begin{itemize}
 \item[$\Box$] 投稿経験が少ない人は，採録された経験の豊富な人に校正して
	       もらう．
 \item[$\Box$] 読者の立場から見て論理的な飛躍がないかに注意して記述する．
\end{itemize}

%5.7
\subsection{その他}

\begin{itemize}
 \item[$\Box$] 投稿前にチェックリストの各項目を満たしているか，必ず確認
	       する． 
\end{itemize}

%6
\section{おわりに}

本稿では，A4縦型2段組み用に変更したスタイルファイルを用いた論文のフォー
マット方法と，論文誌ジャーナル編集委員会がまとめた「べからず集」に基づく
論文の書き方を示した．内容的にまだ不十分の部分が多いため，意見，要望等を
\begin{quote}
 \|editt@ipsj.or.jp|
\end{quote}
までお寄せ頂きたい．
\cite{boutet_being_2016}


\begin{acknowledgment}
A4横型に対するガイドを基に，本稿を作成した．
クラスファイルの作成においては，
京都大学の中島　浩氏にさまざまなご教示を頂き，
さらにBiB\TeX 関連ファイルの利用についても快諾頂いたことを深謝する．
また，A4横型に対するガイドを作成された当時の編集委員会の担当者に深謝する．
\end{acknowledgment}


\bibliographystyle{ipsjunsrt}
\bibliography{references}


\begin{thebibliography}{10}

\bibitem{okumura}
奥村晴彦：改訂第5版\LaTeXe 美文書作成入門，
技術評論社(2010).

\bibitem{companion}
Goossens, M., Mittelbach, F. and Samarin, A.:
{\it The LaTeX Companion},
Addison Wesley, Reading, Massachusetts (1993).

\bibitem{book1}
木下是雄：
理科系の作文技術，
中公新書(1981).

\bibitem{book2}
Strunk W. J. and White E.B.:
{\it The Elements of Style, Forth Edition},
Longman (2000).

\bibitem{book3}
Blake G. and Bly R.W.:
{\it The Elements of Technical Writing},
Longman (1993).

\bibitem{book4}
Higham N.J.:
{\it Handbook of Writing for the Mathematical Sciences},
SIAM (1998).

\bibitem{webpage1}
情報処理学会論文誌ジャーナル編集委員会：
投稿者マニュアル(online)，
\urlj{http://www.ipsj.or.jp/journal /submit/manual/j\_manual.html}
(2007.04.05).

\bibitem{webpage2}
情報処理学会論文誌ジャーナル編集委員会：
べからず集(online)，
\urlj{http://www.ipsj.or.jp/journal/manual /bekarazu.html}
(2011.09.15).

\end{thebibliography}




\end{document}
